\documentclass[11pt,a4paper]{article}
% ============================================================
% Shared preamble for Claude Experiments reports (Ukrainian, LuaLaTeX)
% Compile with: lualatex report_*_uk.tex
% Include this AFTER \documentclass in each report:
%   \documentclass[11pt,a4paper]{article}
%   % ============================================================
% Shared preamble for Claude Experiments reports (Ukrainian, LuaLaTeX)
% Compile with: lualatex report_*_uk.tex
% Include this AFTER \documentclass in each report:
%   \documentclass[11pt,a4paper]{article}
%   % ============================================================
% Shared preamble for Claude Experiments reports (Ukrainian, LuaLaTeX)
% Compile with: lualatex report_*_uk.tex
% Include this AFTER \documentclass in each report:
%   \documentclass[11pt,a4paper]{article}
%   \input{../preamble_uk}
% ============================================================
% Cyrillic via LuaLaTeX + fontspec (pdflatex has no T2A here).
\usepackage{fontspec}
\setmainfont{DejaVu Serif}
\setsansfont{DejaVu Sans}
\setmonofont{DejaVu Sans Mono}
\usepackage[a4paper,top=2.5cm,bottom=2.5cm,left=2.5cm,right=2.5cm]{geometry}
\usepackage{amsmath,amssymb}
\usepackage{graphicx}
\usepackage{subcaption}
\usepackage{booktabs}
\usepackage{xcolor}
\usepackage[colorlinks=true,allcolors=blue]{hyperref}
\usepackage{microtype}
\usepackage{float}
\usepackage{array}
\microtypesetup{expansion=false}

% Color scheme
\definecolor{colMicro}{HTML}{222222}
\definecolor{colExact}{HTML}{2878b5}
\definecolor{colPlacement}{HTML}{e07a1f}
\definecolor{colBoundary}{HTML}{429c74}
\definecolor{colCurrent}{HTML}{c23b55}
\definecolor{colMM}{HTML}{c23b55}

\providecommand{\SIG}{\textcolor{green!50!black}{\textbf{*}}}
\providecommand{\NS}{\textcolor{gray}{n.s.}}
\providecommand{\DRMS}{\Delta\mathrm{RMST}}
\providecommand{\ci}[2]{$[{#1},\;{#2}]$}

% ============================================================
% Cyrillic via LuaLaTeX + fontspec (pdflatex has no T2A here).
\usepackage{fontspec}
\setmainfont{DejaVu Serif}
\setsansfont{DejaVu Sans}
\setmonofont{DejaVu Sans Mono}
\usepackage[a4paper,top=2.5cm,bottom=2.5cm,left=2.5cm,right=2.5cm]{geometry}
\usepackage{amsmath,amssymb}
\usepackage{graphicx}
\usepackage{subcaption}
\usepackage{booktabs}
\usepackage{xcolor}
\usepackage[colorlinks=true,allcolors=blue]{hyperref}
\usepackage{microtype}
\usepackage{float}
\usepackage{array}
\microtypesetup{expansion=false}

% Color scheme
\definecolor{colMicro}{HTML}{222222}
\definecolor{colExact}{HTML}{2878b5}
\definecolor{colPlacement}{HTML}{e07a1f}
\definecolor{colBoundary}{HTML}{429c74}
\definecolor{colCurrent}{HTML}{c23b55}
\definecolor{colMM}{HTML}{c23b55}

\providecommand{\SIG}{\textcolor{green!50!black}{\textbf{*}}}
\providecommand{\NS}{\textcolor{gray}{n.s.}}
\providecommand{\DRMS}{\Delta\mathrm{RMST}}
\providecommand{\ci}[2]{$[{#1},\;{#2}]$}

% ============================================================
% Cyrillic via LuaLaTeX + fontspec (pdflatex has no T2A here).
\usepackage{fontspec}
\setmainfont{DejaVu Serif}
\setsansfont{DejaVu Sans}
\setmonofont{DejaVu Sans Mono}
\usepackage[a4paper,top=2.5cm,bottom=2.5cm,left=2.5cm,right=2.5cm]{geometry}
\usepackage{amsmath,amssymb}
\usepackage{graphicx}
\usepackage{subcaption}
\usepackage{booktabs}
\usepackage{xcolor}
\usepackage[colorlinks=true,allcolors=blue]{hyperref}
\usepackage{microtype}
\usepackage{float}
\usepackage{array}
\microtypesetup{expansion=false}

% Color scheme
\definecolor{colMicro}{HTML}{222222}
\definecolor{colExact}{HTML}{2878b5}
\definecolor{colPlacement}{HTML}{e07a1f}
\definecolor{colBoundary}{HTML}{429c74}
\definecolor{colCurrent}{HTML}{c23b55}
\definecolor{colMM}{HTML}{c23b55}

\providecommand{\SIG}{\textcolor{green!50!black}{\textbf{*}}}
\providecommand{\NS}{\textcolor{gray}{n.s.}}
\providecommand{\DRMS}{\Delta\mathrm{RMST}}
\providecommand{\ci}[2]{$[{#1},\;{#2}]$}


\title{CE-04: Похибка MicroMacro за динаміки SIR\\
       \large Звіт про експеримент --- проєкт Two-Scale CTMC}
\author{Claude Experiments Pipeline}
\date{\today}

\begin{document}
\maketitle

\begin{abstract}
Цей звіт досліджує, як похибка апроксимації MicroMacro масштабується з
інтенсивністю одужання $\gamma$ в епідемічній моделі SIR.
Усі сценарії мають спільні $\beta = 0.003$, топологію chain-of-cliques
($L=10$, $n=100$, $d=6$, $k=5$) та 100 парних реплікацій.
Похибка $\DRMS$ (MicroMacro $-$ повний Gillespie) значуще від'ємна
в усіх сценаріях і монотонно зростає з $\gamma$.
Одужання посилює передчасний export hazard, оскільки macro layer
не враховує проходження хвилі: boundary-вузли можуть уже
одужати на момент спрацювання експортованого macro hazard.
\textit{Ukrainian summary}: див.\ \texttt{docs/ce04\_sir\_validation/report.md}.
\end{abstract}

\section{Дизайн експерименту}

\subsection{Мета}
Визначити, чи додавання SIR-одужання ($\gamma > 0$) до SI baseline
підсилює чи зменшує похибку апроксимації MicroMacro.

\subsection{Мережа та параметри}
\begin{itemize}
  \item Топологія: chain-of-cliques, $L=10$ громад, по $n=100$ вузлів у кожній,
        степінь $d=6$, $k=5$ інтеркомунітарних ребер.
  \item Інтенсивність інфікування: $\beta = 0.003$ (у межах і між громадами).
  \item Інтенсивності одужання: $\gamma \in \{0,\; 0.0005,\; 0.001,\; 0.003\}$.
  \item Базове репродуктивне число: $R_0 = \beta/\gamma \in \{\infty, 6, 3, 1\}$.
  \item Реплікації: 100 парних прогонів (той самий seed для Micro та MicroMacro).
  \item Горизонт: $T = 5000$ одиниць часу; епідемія вважається завершеною,
        коли не залишається інфікованих вузлів або $t > T$.
\end{itemize}

\subsection{Первинна метрика}
\[
  \DRMS = \overline{T_{\text{MM}}} - \overline{T_{\text{Micro}}}
\]
де $T$ позначає середній час першого прибуття до останньої громади.
Від'ємний $\DRMS$ означає передчасне прибуття за MicroMacro.
Bootstrap 95\% довірчі інтервали використовують 10{,}000 ресемплів.

\section{Результати}

\begin{table}[H]
\centering
\caption{Результати CE-04: похибка MicroMacro за інтенсивністю одужання.
  $\DRMS$ = різниця середніх (MicroMacro $-$ Micro).
  Значуще, якщо CI виключає нуль.}
\label{tab:ce04-main}
\begin{tabular}{lrrrrl}
\toprule
Scenario & $\gamma$ & $R_0$ & $\DRMS$ & 95\% CI & Sig. \\
\midrule
SI        & 0.0000 & $\infty$ & $-1067.6$ & $[-1242.5,\; -917.7]$  & \SIG \\
SIR weak  & 0.0005 & 6.0    & $-1244.3$ & $[-1428.3,\; -1054.1]$ & \SIG \\
SIR mid   & 0.0010 & 3.0    & $-1587.3$ & $[-1788.0,\; -1375.2]$ & \SIG \\
SIR strong& 0.0030 & 1.0    & $-1726.3$ & $[-2020.4,\; -1426.1]$ & \SIG \\
\bottomrule
\end{tabular}
\end{table}

\noindent Усі чотири сценарії статистично значущі.
Похибка монотонно зростає з $\gamma$ (рядки 1--4 фінальні).

\subsection{Відносна відстань Wasserstein}

Відносна $W_1$ відстань та частки завершення епідемії:
\begin{table}[H]
\centering
\caption{Частки завершення та похибка розподілів за сценаріями.
  \emph{Completion} = частка зі 100 реплікацій, у яких епідемія
  досягла останньої громади.}
\label{tab:ce04-completion}
\begin{tabular}{lrrrr}
\toprule
Scenario & $\gamma$ & Micro compl. & MM compl. & rel-$W_1$ \\
\midrule
SI        & 0.0000 & 100\% & 100\% & 0.3706 \\
SIR weak  & 0.0005 &  89\% &  96\% & 0.4102 \\
SIR mid   & 0.0010 &  70\% &  93\% & 0.4703 \\
SIR strong& 0.0030 &  28\% &  72\% & 0.4399 \\
\bottomrule
\end{tabular}
\end{table}

\textbf{Критичне спостереження:} при $\gamma = 0.003$ ($R_0 = 1$) лише 28\%
реплікацій повного Gillespie породжують епідемію на весь ланцюг, але MicroMacro
оголошує завершення в 72\% реплікацій --- \textbf{false-positive rate
у 44 процентні пункти}.  Це якісно інший режим збою:
MicroMacro не просто прогнозує раннє прибуття, він прогнозує \emph{прибуття
тоді, коли жодного прибуття не сталося б} у стохастичній реальності.

\section{Інтерпретація}

\subsection{Механізм підсилення похибки}

За двомасштабної апроксимації macro layer спрацьовує import event,
використовуючи агреговану частку $\bar{i}(t)$ інфікованих вузлів на момент
експорту.  У моделі SIR $\bar{i}(t)$ спершу зростає (хвиля інфекції), а
потім спадає (одужання).  Є два режими збою:

\begin{enumerate}
  \item \textbf{Wave-passage mismatch.}  Після піку $\bar{i}(t)$
        спадає.  Якщо macro transition спрацьовує, коли хвиля вже пройшла,
        hazard менший, ніж був на піку хвилі, проте подія
        спрацьовує, бо агрегована macro rate інтегрується за повною
        формою хвилі, а не за поточною.  Результат: штучно ранні прибуття.

  \item \textbf{Censoring effect.}  За великого $\gamma$ багато реплікацій
        досягають $R_0 \approx 1$ і не породжують епідемії на весь ланцюг.
        Ці цензуровані прогони (без прибуття) відсутні в розподілі
        MicroMacro, але роздувають середнє Micro --- або відсутні в обох,
        залежно від реалізації.  \emph{Умовне} середнє за
        MicroMacro тоді значно від'ємніше.
\end{enumerate}

\subsection{Значення для статті}

Випадок SI --- первинний сценарій статті.  Результати CE-04 показують:
\begin{itemize}
  \item MicroMacro передчасний навіть у SI baseline ($\DRMS \approx -1068$).
  \item Одужання ($\gamma > 0$) монотонно \emph{погіршує} цю похибку.
  \item Твердження статті про передчасний export hazard тому є консервативним:
        реальна SIR-похибка більша за наведену SI-похибку.
  \item Поблизу $R_0 = 1$ ($\gamma = 0.003$) MicroMacro генерує
        false-positive rate у 44 процентні пункти для завершення епідемії.
        Це якісно відмінний збій: модель прогнозує
        епідемії, які \emph{ніколи не сталися б} у стохастичній симуляції.
  \item Відстань rel-$W_1$ немонотонна при $\gamma = 0.003$
        (спадає з 0.4703 до 0.4399), бо більшість цензурованих прогонів Micro
        не мають часу прибуття, стискаючи референсний розподіл.
\end{itemize}

\section{Висновок}

Одужання підсилює похибку апроксимації MicroMacro двома різними способами.
По-перше, умовний $\DRMS$ монотонно зростає: з $-1068$ (SI) до
$-1726$ (SIR, $\gamma=0.003$), тобто на 62\%.
По-друге, і більш критично, поблизу епідемічного порогу ($R_0=1$),
MicroMacro генерує на 44 процентні пункти більше хибних завершень епідемії,
ніж повна стохастична модель.
Обидва режими збою простежуються до wave-passage mismatch: агрегований
macro hazard не враховує ані пост-пікового спаду $\bar{i}(t)$,
ані ймовірності того, що окремі boundary-вузли вже одужали.
Майбутня робота має розвинути bridge-aware hazard корекції, що
обумовлені станом інфікованості gateway-вузлів, а не агрегованою громадською часткою.

\section*{Розташування даних}
\texttt{output/claude\_experiments/ce04\_sir\_validation/}

\end{document}
